% !TeX root = ../papers/nc-p3.tex
\section{第二長方形の特別成分での無印商と一次元安定群}
\label{sec:parallel-vertical-orbits}
\begin{theorem}\label{thm:parallel-vertical-orbits}
112型の $b=c=1,t=\pm1$ における正規化した八代表台で、全ての正則点の指定始点の同型と最小周期は次表で尽くされる。
$B=1$, $(a,y,d)=(A/t,C,D/t)$ とし、各行で台と行列式の非零条件および
定理\ref{thm:parallel-support-fibers}の正則条件を課す。$\sigma=a+y$, $\pi=ay$ と書く。
\[
\begin{array}{c|c|c|c}
 台&\text{係数の等式}&\text{完全不変量（各符号で）}&\text{最小周期}\\\hline
118&a=y=1,\ d=0&\text{一点}&2\\
124&a=y=1&d&2\\
129&a=1,\ y=0&d&1\\
134&a=1,\ d=0&y&2\\
139&a=1&(y,d)&2\\
142&\sigma=1,\ d=0&\pi&1\\
145&d=\sigma-1&(\sigma,\pi)&2\\
147&\text{開条件のみ}&(\sigma,\pi,d)&1\ (\sigma=1),\ 2\ (\sigma\ne1).
\end{array}
\]
異なる行、異なる符号の間に指定始点の同型はない。台147の商は例えば
$\pi(1-\sigma+\pi)d(\sigma-d-1)(\pi-d)\ne0$ の開集合である。
$\sigma^2=4\pi$ 上では $t=1$ の射影安定群が位数8の二面体群になり、$t=-1$ では
$\mathbb G_m\rtimes\mathbb Z/2$（逆数作用）になる。$a\ne y$ では両符号で標準の四元群である。
\end{theorem}

\begin{proof}[固有空間の分割の有限検査]
各位置で基底に関して対角な内在的作用素の同時固有空間分解を作る。
まず各 $p\ne q$ について $W_{pq}G_{pq}^t$ の平方の部分traceを使う。
さらに他の三位置の既知の固有空間を一つずつ選び、その射影を挿入した $W,G$ の縮約を加える。
位置 $p$ の文字 $i$ での最後の対角値は、残る三文字が指定の三ブロックに入る語での
$W_kG_k$ の和である。固有値の等しい文字を一組のブロックとし、分割が細かくなる限り繰り返す。
各操作はテンソルの基底変更に関して自然で、相対不変量の共通倍率は分割に影響しない。

二位置回転は偶数位置で恒等、奇数位置で四元群の1を加える置換と対角倍率で元の窓に戻る。
従って最終分割は位置0,1の分割 $A,B$ から $(A,B,A,B\oplus1)$ と書ける。
$a\ne y$ では目標は四つの一元集合である。$a=y$ では後述の対称性で偶数位置の $\{0,1\}$、
奇数位置の $\{2,3\}$ は同じブロックに入り、残り二文字を一つずつ分離するのが目標である。
この目標より粗い全候補を取り、そのブロック内の隣り合う二文字について、上の部分traceの差と
全三ブロックへの射影を挿入した差を同時に零とする。これは安定分割であるための必要条件である。
開条件の積、行列式、$a\ne y$ の層では $a-y$ を逆にした環で、これらの差のイデアルは
両符号・全台・全候補で単位イデアルになる。初期の二分割の候補数は非対角で $15^2-1=224$、
対角で $5^2-1=24$ に過ぎない。検算は単元因子を除いてGröbner基底を計算し、
必要な場合は定義域の積で飽和する。\path{computations/verify_rectangle_partition_frames.py} が全条件を生成する。
\end{proof}

\begin{proof}[対角の二文字と二つの安定群]
$a=y$ で四つ組の隣接位置 $0,1$ にそれぞれ二次元固有空間を残し、位置2,3では一元の固有線に射影する。
得られる二行二列行列は $t=1$ なら恒等と交換行列の張る鉛筆で、$t=-1$ ならこれに固定した符号行列を掛けた鉛筆となる。
位置2,3に二次元空間を残しても同様である。いずれも階数一の二点から各二次元空間の和・差の二直線が復元できる。
二次元空間内で和と差を基底とした表示では全同型が単項となり、元の一元固有線の対と復元した二線の対は区別される。

この新表示で指定始点の全台置換は恒等および全位置の和・差の同時交換のみである。
後者は元の四元群の文字符号ですでに実現される。各語の係数を定数と既約因子の積で表示し、
対角倍率の重み行列の左整数核を係数の指数に作用させる。得た格子のHermite標準形は、
正次元の各層で全て単位行列である。例えば台147での因子は
$a,d,a-1,2a-d-1,a^2-d$ であり、台145では $a,a-1,2a-1$ になる。
従って係数を一意に復元できる。台および符号が異なる場合の新表示の台置換は存在しない。

$t=1$ で新表示の重み行列の非零Smith因子は $1^{11},2,2$。対角群の位数4と同時交換から
安定群の位数は8となる。元の表示の標準四元群と二文字交換は非可換で、位数8の非可換群で
四元群を含むものは二面体群である。新表示の対角四元群は元のものと異なり、各位置で一つの固有文字が重なる。
$t=-1$ ではSmith因子は $1^{12}$ で、位置スカラーを除く対角部分は一次元トーラスとなる。
二次元空間の和・差の二線の重みは、各位置のスカラーを除くと $1,-1$ または $-1,1$ で、
元の一元線の重みは両方0である。元の文字符号による同時交換はこのトーラスを反転する。
この場合の全四元部分群は、一つのトーラス内対合と外の対合で生成され、トーラス共役で互いに移る。
\end{proof}

\begin{proof}[指定始点と移動]
$a\ne y$ で元の四枠を復元した場合、台147を保つ二置換は恒等と $I(a,y,d)=(y,a,d)$ である。
薄い係数と $B=1$ の規格化で同型を比較できる。一位置移動の代表は
\[
 J(a,y,d)=(1-a,1-y,1-a-y+d).
\]
$J^2=1$, $IJ=JI$ である。各小さい台での台置換は全て台147の対応する置換の制限である。
異なる台への置換はなく、指定始点で $I$ が残るのは118,124,142,145,147のみである。
二位置回転は常に許され、奇数回転を許すのは129,142,147である。
その場合に同じ類に戻る条件は $\sigma=1$ で、このとき $IJ$ が係数を固定する。
$a=y$ の新表示でも奇数回転が可能な対角の台は142,147だけで、上の変換が実現する同型と全置換数が一致する。
非対角と対角は復元した分割の次元で区別される。
元の枠が四線に分かれる各台の対角Smith因子は $1^{11},2,2$ であり、$a\ne y$ では恒等以外の置換は係数を固定しない。
\path{computations/verify_parallel_vertical_orbits.py} が薄い規格化、全軌道の台置換、和・差の行列鉛筆、両種の格子を検算する。
\end{proof}
